% !TeX program = xelatex
% Encoding: UTF-8. Compile with XeLaTeX; BibTeX/Biber are not required.
% Tables and figures use saved output from mid_term_code.ipynb.
\documentclass[12pt,a4paper]{article}

\usepackage{iftex}
\ifPDFTeX
  \PackageError{bao-cao-midterm}{Select XeLaTeX or LuaLaTeX}{This document uses Unicode fonts.}
\fi
\usepackage{fontspec}
\defaultfontfeatures{Ligatures=TeX}
% Bundled font with Vietnamese glyphs for names and university details on the cover.
\newfontfamily\coverfont{LiberationSerif}[Path=fonts/,Extension=.ttf,
  UprightFont=*-Regular,BoldFont=*-Bold]
% The source archive includes Times-compatible OpenType fonts for a stable layout.
\IfFileExists{fonts/NimbusRoman-Regular.otf}{%
  \setmainfont{NimbusRoman}[Path=fonts/,Extension=.otf,UprightFont=*-Regular,
    BoldFont=*-Bold,ItalicFont=*-Italic,BoldItalicFont=*-BoldItalic]}{%
  \IfFontExistsTF{Times New Roman}{\setmainfont{Times New Roman}}{\setmainfont{TeX Gyre Termes}}}
\IfFileExists{fonts/NimbusSans-Regular.otf}{%
  \setsansfont{NimbusSans}[Path=fonts/,Extension=.otf,UprightFont=*-Regular,
    BoldFont=*-Bold,ItalicFont=*-Italic,BoldItalicFont=*-BoldItalic]}{\setsansfont{TeX Gyre Heros}}
\IfFontExistsTF{TeX Gyre Cursor}{\setmonofont{TeX Gyre Cursor}[Scale=MatchLowercase]}{%
  \setmonofont{DejaVu Sans Mono}[Scale=MatchLowercase]}
\usepackage[english]{babel}
\usepackage[left=3cm,right=2cm,top=2.5cm,bottom=2.5cm,headheight=16pt]{geometry}
\usepackage{amsmath,amssymb}
\usepackage{graphicx}
\graphicspath{{figures/}{report/figures/}}
\usepackage{booktabs,tabularx,longtable,array,makecell}
\usepackage{caption,subcaption}
\usepackage{placeins,float}
\usepackage{titlesec,needspace}
\usepackage{enumitem}
\usepackage{xcolor}
\definecolor{navy}{HTML}{17334C}
\definecolor{teal}{HTML}{087F8C}
\definecolor{muted}{HTML}{526474}
\definecolor{light}{HTML}{EDF4F8}
\usepackage{setspace}
\usepackage{fancyhdr}
\usepackage{xurl}
\usepackage[unicode,colorlinks=true,linkcolor=navy,citecolor=teal,urlcolor=teal]{hyperref}
\usepackage{bookmark}
\usepackage{microtype}

\hypersetup{
  pdftitle={Midterm Project Report: Book Rating Classification},
  pdfsubject={Assignment 6 --- Classification: Book Recommendation Dataset},
  pdfkeywords={Python, data analysis, Book-Crossing, classification, target encoding}
}
\setstretch{1.20}
\raggedbottom
\clubpenalty=10000
\widowpenalty=10000
\displaywidowpenalty=10000
\setcounter{topnumber}{2}
\setcounter{bottomnumber}{2}
\setcounter{totalnumber}{4}
\renewcommand{\topfraction}{0.85}
\renewcommand{\bottomfraction}{0.75}
\renewcommand{\textfraction}{0.10}
\renewcommand{\floatpagefraction}{0.70}
\titleformat{\section}{\sffamily\bfseries\fontsize{15}{18}\selectfont\color{navy}}{\thesection}{0.7em}{}
\titleformat{\subsection}{\sffamily\bfseries\fontsize{12.5}{15}\selectfont\color{navy}}{\thesubsection}{0.7em}{}
\titlespacing*{\section}{0pt}{1.25\baselineskip}{0.45\baselineskip}
\titlespacing*{\subsection}{0pt}{0.95\baselineskip}{0.30\baselineskip}
\setlength{\textfloatsep}{13pt plus 3pt minus 2pt}
\setlength{\intextsep}{12pt plus 3pt minus 2pt}
\renewcommand{\arraystretch}{1.12}
\setlength{\parindent}{1.2em}
\setlength{\parskip}{0.3em}
\setlength{\emergencystretch}{3em}
\setcounter{tocdepth}{2}
\captionsetup{font=footnotesize,labelfont=bf,justification=raggedright,singlelinecheck=false,skip=7pt}
\setlist[itemize]{topsep=3pt,itemsep=2pt,leftmargin=1.6em}
\setlist[enumerate]{topsep=3pt,itemsep=2pt,leftmargin=1.8em}
\newcolumntype{Y}{>{\raggedright\arraybackslash}X}
\newcolumntype{L}[1]{>{\raggedright\arraybackslash}p{#1}}
\newcommand{\code}[1]{\texttt{#1}}
\newcommand{\source}[1]{\par\smallskip{\raggedright\footnotesize\color{muted}Source: #1.\par}}
\newcommand{\reportfigure}[5][1.00]{%
  \begin{figure}[htbp]
    \centering
\includegraphics[width=#1\linewidth,height=0.36\textheight,keepaspectratio]{#2}
    \caption{#3}\label{#4}
    \source{#5}
  \end{figure}%
}
\pagestyle{fancy}
\fancyhf{}
\fancyhead[L]{\small\sffamily\color{muted}BOOK RATING CLASSIFICATION}
\fancyhead[R]{\small\sffamily\color{muted}ASSIGNMENT 06}
\fancyfoot[L]{\footnotesize\color{muted}Data Analysis with Python}
\fancyfoot[R]{\footnotesize\thepage}
\renewcommand{\headrulewidth}{0.3pt}
\renewcommand{\footrulewidth}{0pt}

\begin{document}
\hypersetup{pageanchor=false}
\begin{titlepage}
  \thispagestyle{empty}
  \begin{center}
  \includegraphics[width=4.2cm]{neu_logo.png}\par
  \vspace{0.2cm}
  {\coverfont\bfseries\large\color{navy}ĐẠI HỌC KINH TẾ QUỐC DÂN\par}
  \vspace{0.6cm}
  {\sffamily\large\color{teal}MIDTERM PROJECT REPORT\par}
  \vspace{0.25cm}
  {\noindent\color{teal}\rule{\linewidth}{1.5pt}\par}
  \vspace{0.4cm}
  {\sffamily\bfseries\Huge\color{navy}BOOK RATING\\[0.15cm]CLASSIFICATION\par}
  \vspace{0.35cm}
  {\sffamily\Large Book Recommendation Dataset\par}
  \vspace{0.2cm}
  {\large Assignment 6 --- Classification\par}
  \end{center}
  \vspace{0.35cm}
  \noindent\colorbox{light}{%
    \parbox{0.93\linewidth}{\vspace{0.25cm}
      \textbf{Research Question}\\[0.15cm]
      Can a rating be classified as high (8--10) or not high (1--7) using user and book
      information together with other ratings in the training set?
      \vspace{0.25cm}}
  }
  \vspace{0.35cm}
  \begin{center}
  \begin{tabular}{ccc}
    \textbf{433,671} & \textbf{0.6553} & \textbf{0.7306} \\
    Explicit ratings & Test Macro F1 & Test ROC-AUC
  \end{tabular}
  \end{center}
  \vspace{0.25cm}
  {\coverfont
  \noindent\textbf{Thành viên thực hiện:}\par
  \noindent Lương Văn Thái Tuấn -- 11258665\\
  Nguyễn Thị Khánh Hòa -- 11253156\par
  \medskip
  \noindent\textbf{Giảng viên hỗ trợ:} Tiến sĩ Trần Đức Minh\par}
  \vfill
  \noindent\textbf{Notebook:} \nolinkurl{mid_term_code.ipynb}\\
  \textbf{Course:} Data Analysis with Python\\
  \textbf{Selected model:} Logistic Regression + Target Encoding\\
  \textbf{Repository:} \url{https://github.com/Bluenip/Data_analyst_python_midterm_project}\\
  \textbf{Report date:} 2 November 2026
\end{titlepage}
\hypersetup{pageanchor=true}
\pagenumbering{roman}

\Needspace{10\baselineskip}
\section*{Abstract}
\addcontentsline{toc}{section}{Abstract}
This report uses the Books, Users, and Ratings tables from the Book Recommendation Dataset to
formulate a binary classification task at the individual-rating level. Only ratings from 1 to 10
are used as labels; ratings of 8 or higher are designated class 1. Zero-valued records are retained
when describing rating activity but are not treated as low ratings. After integration and cleaning,
the classification dataset contains \textbf{433,671 ratings}, of which \textbf{57.62\%} belong to class 1.

The workflow covers key validation and encoding checks, ISBN normalization, missing and anomalous
value handling, exploratory data analysis (EDA), feature construction, a stratified 70/15/15
train/validation/test split, and comparisons of four algorithms across five feature groups. Each
algorithm uses the same configuration across feature groups. Imputation, scaling, and one-hot
encoding are fit on training data only; user- and book-level rates are generated with smoothed
out-of-fold (OOF) encoding to avoid using a row's own label to encode that row.

Of the 20 evaluated configurations, Logistic Regression with Target Encoding was selected by
validation Macro F1. On the test set, it achieved an Accuracy of \textbf{65.65\%}, a Macro F1 of
\textbf{0.6553}, and a ROC-AUC of \textbf{0.7306}. Compared with a baseline that always predicts the
majority class, Macro F1 increased by \textbf{0.2897} and Accuracy by \textbf{8.03 percentage points}.
The results indicate that rating-history features from the training set provide predictive signal,
but performance declines for cases with both an unseen user and an unseen book. The evaluation covers one random rating-level split, with a test set already inspected during
earlier development. Results therefore describe this experiment rather than an independent
confirmation of future performance or recommendation ranking quality.

\clearpage
\begingroup
\footnotesize
\setstretch{1.0}
\setlength{\parskip}{0pt}
\tableofcontents
\endgroup
\clearpage
\listoffigures
\listoftables
\clearpage
\pagenumbering{arabic}

\Needspace{10\baselineskip}
\section{Problem Definition and Scope}
\Needspace{5\baselineskip}
\subsection{Objective}
The unit of observation is one user's rating of one book. The objective is to estimate whether a
rating belongs to the high-rating range. The task is framed as binary classification rather than
exact score prediction:
\begin{equation}
 y_i = \operatorname{liked}_i =
 \begin{cases}
  1, & \text{if } 8 \le r_i \le 10,\\
  0, & \text{if } 1 \le r_i \le 7.
 \end{cases}
 \label{eq:target}
\end{equation}
Here, $r_i$ denotes the \code{Book-Rating} for observation $i$. The threshold of 8 is a project
convention, not a requirement of the midterm assignment. Class 0 means only that a rating does not
meet the high-rating threshold; it does not imply that the user disliked the book.

According to the data provider, ratings from 1 to 10 are explicit, while 0 represents implicit
feedback~\cite{dataset,description}. Because zero-valued records do not provide an explicit
preference score, they are not assigned to either class in Equation~\eqref{eq:target}.

\Needspace{5\baselineskip}
\subsection{Research Questions}
\begin{enumerate}
  \item What patterns characterize the class distribution, user ages, and user activity levels?
  \item Do book metadata, country, and ID-based statistics improve validation classification
  performance over the initial age- and year-based features?
  \item Which configuration best identifies both classes and improves on the baseline?
  \item Does aggregate performance reflect results for users or books unseen during training?
\end{enumerate}

\Needspace{5\baselineskip}
\subsection{Interpretation Scope}
The data represent only ratings present in the collected dataset. The model classifies a rating
known to have a score from 1 to 10; it does not predict the probability that an arbitrary user will
read, purchase, or rate a book. Classification metrics also do not replace ranking metrics for a
recommendation system. Findings from EDA are descriptive or hypothesis-generating. The main sequence is data validation
and cleaning, EDA, feature construction, validation-based selection, and test error analysis.

\Needspace{10\baselineskip}
\section{Data, Validation, and Integration}
\Needspace{5\baselineskip}
\subsection{Source and Structure}
The data come from Möbius's Book Recommendation Dataset on Kaggle, the source specified for
Assignment 6~\cite{dataset}. The three tables provide distinct types of information, summarized in
Table~\ref{tab:raw}.
\begin{table}[htbp]
\centering
\caption{Size of the three raw data tables.}\label{tab:raw}
\begin{tabular}{lrr}
\toprule
Table & Rows & Columns \\
\midrule
Books & 271,360 & 8 \\
Users & 278,858 & 3 \\
Ratings & 1,149,780 & 3 \\
\bottomrule
\end{tabular}
\source{Notebook, Section 1: Data Inspection}
\end{table}

\begin{itemize}
  \item \textbf{Books:} ISBN, title, author, publication year, publisher, and three cover-image URLs.
  \item \textbf{Users:} User-ID, Location, and Age.
  \item \textbf{Ratings:} User-ID, ISBN, and Book-Rating.
\end{itemize}
Image URLs are excluded because the project does not process images. ISBNs are loaded as strings to
preserve leading zeros and the character X. All three files are checked using strict UTF-8 decoding
before being read with \code{pandas.read\_csv}; a loop automates the loading process.

\Needspace{5\baselineskip}
\subsection{Key and Initial Quality Checks}
In the raw data, ISBNs in Books and User-IDs in Users are unique. There are no exact duplicate rows
or duplicate User-ID--ISBN pairs in Ratings. A total of 118,644 ratings have ISBNs that do not match
Books; every User-ID in Ratings matches Users. The checks also find one Books ISBN and 14 Ratings
ISBNs with leading or trailing whitespace. These join-key issues must be resolved before creating
the clean analysis table, even though the original keys are unique.

Raw ages range from 0 to 244 and require review before use. Three book rows contain publisher names
in the publication-year field, with some following fields containing misaligned URLs. These rows are
isolated during cleaning rather than corrected by guesswork.

\Needspace{5\baselineskip}
\subsection{Relationships and Join Strategy}
Ratings is the central table. Each user may contribute multiple ratings, and each book may appear in
multiple ratings, so both relationships are many-to-one from Ratings to their respective lookup tables:
\begin{center}
\begin{tabular}{c}
\code{Ratings.User-ID} $\longrightarrow$ \code{Users.User-ID}\\[2pt]
\code{Ratings.ISBN} $\longrightarrow$ \code{Books.ISBN}.
\end{tabular}
\end{center}
The project uses a \textbf{left join}, not concatenation, because the three tables have different
structures. The \code{validate='many\_to\_one'} parameter ensures that duplicate keys in a lookup
table do not multiply rows; \code{indicator} records match status. After each join, the row count is
checked against the number of Ratings records. Retaining records without book metadata avoids
silently shrinking the sample; missing values are handled explicitly during preprocessing.
\FloatBarrier

\Needspace{10\baselineskip}
\section{Cleaning and Classification Dataset Construction}
\Needspace{5\baselineskip}
\subsection{Cleaning Rules}
\begin{enumerate}
  \item \textbf{ISBN and duplicates:} strip leading and trailing whitespace from ISBNs in Books and
  Ratings, then remove exact duplicate rows. One duplicate Books row is removed after normalization.
  If conflicting metadata remain for an ISBN or conflicting ratings for a User-ID--ISBN pair, the
  program stops for review rather than selecting a record automatically.
  \item \textbf{Publication year:} convert values to numeric and isolate the three rows that cannot
  be converted because of field misalignment. Year 0 and years outside 1450--2026 are treated as
  missing. This is a project cleaning rule, not a claim that every value outside this range is invalid.
  \item \textbf{Age:} retain ages from 5 to 100 and set other values to missing without deleting users.
  The number of users with missing age increases from 110,762 to 112,010 after this step.
  \item \textbf{Strings:} normalize whitespace, decode HTML entities in book titles, and replace
  missing or empty author, publisher, and location values in the lookup tables with \code{Unknown}.
  Missing values introduced by the left join are retained and handled in the pipeline.
  \item \textbf{Missingness indicators:} create \code{Age\_missing} and \code{Year\_missing}, then
  update them after joining to include books with no matching metadata. Median imputation is not
  performed at this stage.
  \item \textbf{Ratings and target:} validate integer ratings in the range 0--10 and store them as a
  small integer type; use only ratings from 1 to 10 to create \code{liked}, as defined in
  Equation~\eqref{eq:target}.
\end{enumerate}

The clean joined table retains all 1,149,780 ratings. After key normalization and isolation of
malformed rows, 118,645 ratings have no clean book metadata. This total is one higher than the raw
unmatched-ISBN count because normalization and quarantining change which book keys are available. Within the 433,671-rating classification
dataset, 49,829 records have no book metadata, or 11.49\%. These counts use different denominators
and should not be interchanged.

\Needspace{5\baselineskip}
\subsection{Remaining Missing Values and Training-Time Handling}
\begin{table}[htbp]
\centering
\caption{Missing data in the classification dataset after the clean join.}\label{tab:missing}
\begin{tabularx}{\linewidth}{Yrr}
\toprule
Field & Missing ratings & Percentage \\
\midrule
Age & 130,571 & 30.11\% \\
Year-Of-Publication & 55,635 & 12.83\% \\
Book-Title & 49,829 & 11.49\% \\
Book-Author & 49,829 & 11.49\% \\
Publisher & 49,829 & 11.49\% \\
\bottomrule
\end{tabularx}
\source{Notebook, missing-value table in Section 3; denominator: 433,671 ratings}
\end{table}
Age is missing for 30.11\% of classification ratings but for 40.17\% of users in Users; the
difference arises because one user may contribute multiple ratings. For numeric variables, the
pipeline learns the median from training data for imputation and then scales the values. Missing
categorical values are represented as \code{Unknown}. Missingness indicators allow the model to
distinguish observed values from imputed ones.
\FloatBarrier

\Needspace{10\baselineskip}
\section{Exploratory Data Analysis}
\Needspace{5\baselineskip}
\subsection{Age and Class Distributions}
After cleaning, 166,848 users have valid ages and 112,010 have missing ages. The mean known age
is approximately 34.75 and the median is 32; the 25th and 75th percentiles are 24 and 44.
Figure~\ref{fig:age} summarizes users rather than rating records: each user contributes once.
Most known ages are among younger and middle-aged adults; users with missing ages are excluded.
\reportfigure{01_phan_phoi_tuoi_en.pdf}{User age distribution after applying the 5--100 validity rule; 166,848 users with nonmissing ages.}{fig:age}{Notebook, Section 4.1; chart regenerated from Users}

\begin{samepage}
Class 1 contains \textbf{249,887 ratings} (\textbf{57.62\%}), while class 0 contains
\textbf{183,784 ratings} (\textbf{42.38\%}); see Figure~\ref{fig:labels}.
This is a modest class imbalance. Always predicting class 1 already yields about 57.62\% Accuracy,
so the evaluation also uses Macro F1 and per-class Recall to assess recognition of both groups.
\par\end{samepage}
\reportfigure[0.92]{02_phan_bo_nhan_en.pdf}{Class distribution among the 433,671 explicit ratings; each rating contributes once.}{fig:labels}{Notebook, Section 4.2; chart regenerated from Ratings}
\FloatBarrier

\Needspace{5\baselineskip}
\subsection{Rating Scores and Book Popularity}
Of the 1,149,780 raw ratings, 716,109 have a score of 0 and 433,671 have scores from 1 to 10.
Among explicit ratings, score 8 is most frequent with 103,736 ratings, followed by score 10 with
78,610 and score 7 with 76,457. The distribution is tilted toward relatively high scores, consistent
with class 1 representing more than half of the classification dataset.
\reportfigure[0.92]{03_phan_phoi_diem_en.pdf}{Counts of explicit ratings by score (1--10); zero-valued ratings are excluded.}{fig:ratings}{Notebook, Section 4.2}

\begin{table}[htbp]
\centering
\caption{Five most frequently rated books in Ratings, including zero-valued ratings.}\label{tab:popular}
\small
\begin{tabularx}{\linewidth}{lYr}
\toprule
ISBN & Title & Ratings \\
\midrule
0971880107 & Wild Animus & 2,502 \\
0316666343 & The Lovely Bones: A Novel & 1,295 \\
0385504209 & The Da Vinci Code & 883 \\
0060928336 & Divine Secrets of the Ya-Ya Sisterhood: A Novel & 732 \\
0312195516 & The Red Tent (Bestselling Backlist) & 723 \\
\bottomrule
\end{tabularx}
\source{Notebook, Section 4.4}
\end{table}

Table~\ref{tab:popular} measures rating frequency, not preference. For example, \emph{Wild Animus}
ranks first with 2,502 ratings, but its mean explicit score is only about 4.39 across 581 rated
records. Conversely, the book with ISBN 1888054557 has a mean score of 10.00 but only 11 explicit
ratings. This illustrates why rating counts and mean scores should be considered together; it does
not establish a negative correlation between popularity and scores across all books.

When ranking mean scores, the notebook includes only books with at least 10 explicit ratings to
reduce the influence of books with very few scores. This threshold does not eliminate small-sample
variation; rankings based on 10--20 ratings should be interpreted cautiously.
\FloatBarrier

\Needspace{5\baselineskip}
\subsection{Age and Rating Class}
Among ratings with a known user age, mean ages in classes 0 and 1 are approximately 35.99 and
36.78, respectively. The difference is descriptive; a comparison of means alone cannot establish
how well age classifies individual ratings. When age is grouped into
bins, the proportion of high ratings varies as shown in Figure~\ref{fig:age-target} and
Table~\ref{tab:age-groups}.
\reportfigure{04_tuoi_va_diem_cao_en.pdf}{Share of high ratings (8--10) by age group.}{fig:age-target}{Notebook, Section 4.6}

\begin{table}[htbp]
\centering
\caption{Ratings and class 1 share by age group, restricted to valid ages.}\label{tab:age-groups}
\small
\begin{tabular}{lrr}
\toprule
Age group & Ratings & High-rating share \\
\midrule
5--14 & 3,875 & 60.28\% \\
15--24 & 42,987 & 58.85\% \\
25--34 & 105,871 & 58.57\% \\
35--44 & 73,988 & 59.55\% \\
45--54 & 51,458 & 64.93\% \\
55--64 & 18,963 & 60.45\% \\
65--74 & 4,774 & 59.61\% \\
75--84 & 859 & 53.20\% \\
85--94 & 75 & 34.67\% \\
95--100 & 250 & 76.00\% \\
\bottomrule
\end{tabular}
\source{Notebook, Section 4.6}
\end{table}

The 45--54 age group has 51,458 ratings and a 64.93\% high-rating share. The 85--94 group has a
34.67\% share based on only 75 ratings, while the 95--100 group reaches 76.00\% across 250 ratings.
Extremes in small groups should not be treated as stable patterns. Each row is a rating, so highly
active users contribute more weight to these proportions.

For the $10\times2$ age-group/class table, the notebook reports
$\chi^2(9)=706.6158$ and $p\approx2.61537\times10^{-146}$. As a descriptive effect-size summary,
Cram\'er's $V=\sqrt{\chi^2/(N\min(9,1))}\approx0.0483$, with $N=303{,}100$ ratings,
indicates a small association in this contingency table despite the tiny nominal $p$-value.
The usual independence assumption may fail because users and books contribute repeated ratings;
the $p$-value is therefore a reference calculation, not population-level confirmation. Age remains
a candidate feature to be evaluated through validation performance.
\FloatBarrier

\Needspace{5\baselineskip}
\subsection{User Activity and Rating Tendencies}
The most active user by explicit ratings is ID 11676 with 8,524 ratings, followed by ID 98391 with
5,802. Most users have far fewer ratings. Figure~\ref{fig:age-activity} shows users with a valid known age and at least one explicit rating.
Its logarithmic count axis makes the low-activity majority visible alongside highly active users;
users with missing ages do not appear in this plot.
\reportfigure{05_tuoi_va_hoat_dong_en.pdf}{User age versus explicit rating count, with a logarithmic rating-count axis.}{fig:age-activity}{Notebook, Section 4.7}

To avoid treating all ratings as independent users, the next analysis aggregates each user's class 1
share. The low-activity group contains users at or below the 25th percentile of rating counts
($\leq 1$ rating), while the high-activity group contains users at or above the 75th percentile
($\geq 3$ ratings). The two groups do not overlap in this dataset.
\begin{table}[htbp]
\centering
\caption{Mean user-level high-rating share by activity group.}\label{tab:activity}
\small
\begin{tabularx}{\linewidth}{Yrrr}
\toprule
Group & Users & \makecell{Median\\ratings per user} & \makecell{Mean user-level\\high-rating share} \\
\midrule
Low activity & 45,382 & 1 & 56.27\% \\
High activity & 22,568 & 6 & 53.00\% \\
\bottomrule
\end{tabularx}
\source{Notebook, Section 4.8}
\end{table}
\reportfigure{06_ty_le_diem_cao_theo_hoat_dong_en.pdf}{Comparison of mean user-level high-rating shares across activity groups.}{fig:activity-rate}{Notebook, Section 4.8}
\FloatBarrier

Each user contributes one share to the group mean, regardless of their number of ratings;
users with exactly two ratings belong to neither comparison group. The high-activity group has a
mean high-rating share about 3.28 percentage points lower, using unrounded rates. The
Mann--Whitney U statistic is $U=566,384,643$, with $p\approx3.21475\times10^{-125}$. However, a
user with one rating always has a share of 0 or 1, while a user with many ratings can have a share
between these extremes. Different rating counts therefore produce different variability, so a
difference in distributions does not automatically demonstrate a difference in underlying
preference tendencies. Users may also rate many of the same books.

These results motivate testing activity counts and ID-based statistics. Full-sample EDA summaries
are not used directly as model inputs; features for modeling are recomputed from the training set.
\FloatBarrier

\Needspace{10\baselineskip}
\section{Feature Engineering and Data-Split Design}
\Needspace{5\baselineskip}
\subsection{Five Cumulative Feature Groups}
Table~\ref{tab:features} presents the experimental sequence. Column counts refer to inputs before
one-hot encoding, not the number of dimensions after encoding. Each step adds only the listed group,
while keeping each algorithm's configuration fixed.
\begin{table}[htbp]
\centering
\caption{Input feature groups compared in the experiments.}\label{tab:features}
\small
\begin{tabularx}{\linewidth}{lYc}
\toprule
Group & Features & Columns \\
\midrule
Initial & Age, Age\_missing, Year-Of-Publication, Year\_missing & 4 \\
Metadata & Initial + Book-Author, Publisher & 6 \\
Country & Metadata + country parsed from Location & 7 \\
Count & Country + User\_Rating\_Count, Book\_Rating\_Count & 9 \\
Target Encoding & Count + User\_Liked\_Rate, Book\_Liked\_Rate & 11 \\
\bottomrule
\end{tabularx}
\end{table}

\code{Book-Rating} is excluded from X because the target is derived directly from that column.
User-ID and ISBN are used only as keys for computing statistics, not as continuous numeric inputs.
Book titles, full Location strings, and image URLs are not used directly. EDA does not establish
the independent effects of publication year, author, or publisher; these variables are treated as
candidates and evaluated through validation.

Country is parsed from the final comma-separated component of Location, then normalized to lowercase
and trimmed. Missing markers such as empty strings, \code{n/a}, \code{none}, and \code{unknown} are
mapped to \code{Unknown}; aliases such as \code{us}/\code{united states} are mapped to \code{usa},
and \code{uk} to \code{united kingdom}. Country is unknown for 20,551 ratings, or 4.74\% of the
classification dataset. This rule does not validate all geographic values or correct every
self-reported error.

\Needspace{5\baselineskip}
\subsection{Train, Validation, and Test Split}
Data are randomly split by rating, stratified by label, with \code{random\_state=56}. First, 15\%
is held out for test; then $15/85$ of the remaining data is used for validation.
\begin{table}[htbp]
\centering
\caption{Sizes of the stratified data splits.}\label{tab:split}
\begin{tabular}{lrr}
\toprule
Split & Ratings & Class 1 share \\
\midrule
Train & 303,569 & 57.62\% \\
Validation & 65,051 & 57.62\% \\
Test & 65,051 & 57.62\% \\
\bottomrule
\end{tabular}
\source{Notebook, Section 5.2}
\end{table}
\FloatBarrier

Assertions verify that X and y have matching indices, that the three splits do not overlap, and
that their union reconstructs the classification dataset. Imputation medians, scaling parameters,
and infrequent-category groupings are fit on training data only in the pipeline, following
recommended practices to prevent preprocessing leakage~\cite{pitfalls}. Categorical variables use
one-hot encoding with \code{min\_frequency=100}. Infrequent categories are grouped; a new value
maps to that group if it exists for the feature, and otherwise produces an all-zero one-hot block.

\Needspace{5\baselineskip}
\subsection{Count Features and Target Encoding Without Self-Label Leakage}
The two count features use the number of \emph{ratings scored from 1 to 10 in the training set}.
For a training row, one is subtracted from that ID's count to exclude the current rating. For
validation and test, counts are looked up from training data only; unseen IDs receive 0. These
values are not total book reads or total activity including zero-valued ratings.

For an ID $e$ in fit subset $G$, the smoothed class rate is calculated as:
\begin{equation}
 \operatorname{TE}_{G}(e)=\frac{s_G(e)+\alpha\mu_G}{n_G(e)+\alpha},\qquad
 \mu_G=\frac{1}{|G|}\sum_{j\in G}y_j,\quad\alpha=20,
 \label{eq:te}
\end{equation}
where $n_G(e)$ is the number of ratings and $s_G(e)$ is the sum of class 1 labels for ID $e$ in
subset $G$. The training set is divided into five stratified folds with seed 42; each fold receives
rates computed from the other four folds only. Neither the ID statistics nor the overall rate
$\mu_G$ uses labels from the fold being encoded. For validation and test, $G$ is the full training
set; unseen IDs receive the training-set prior.

Smoothing limits extreme rates for IDs with few ratings, while OOF prevents each row from carrying
its own label information. Here, OOF is a feature-generation step rather than five-fold
cross-validation for the entire model~\cite{encoding}. The number of ratings used to compute training rates differs
from that used for validation/test; this distributional difference remains a limitation when
interpreting results.
\FloatBarrier

\Needspace{10\baselineskip}
\section{Algorithms, Baseline, and Evaluation Criteria}
\Needspace{5\baselineskip}
\subsection{Model Configurations}
Table~\ref{tab:models} lists the four classification algorithms evaluated in the notebook.
Each pipeline has a separately cloned preprocessor and estimator, preventing fitted state from a
previous experiment from being overwritten. The algorithm seed is 42.
\begin{table}[htbp]
\centering
\caption{Configurations held fixed across the five feature groups.}\label{tab:models}
\small
\begin{tabularx}{\linewidth}{lY}
\toprule
Algorithm & Main configuration \\
\midrule
Logistic Regression (LR) & class\_weight=balanced, max\_iter=1000 \\
Decision Tree (DT) & class\_weight=balanced, max\_depth=20, min\_samples\_leaf=50 \\
Linear SVM (SVM) & LinearSVC, class\_weight=balanced, max\_iter=3000 \\
Neural Network (MLP) & Two hidden layers (64, 32), ReLU, Adam; alpha=0.001; batch=1024;
learning\_rate\_init=0.001; max. 100 epochs; tol=0.001; n\_iter\_no\_change=5;
early\_stopping=False \\
\bottomrule
\end{tabularx}
\end{table}

LR and Linear SVM provide linear classifiers; DT can represent nonlinear splits through
condition-based rules; and MLP learns nonlinear transformations through hidden layers. The project
does not perform exhaustive hyperparameter optimization. The three algorithms that support
\code{class\_weight} use the balanced setting; MLP retains default sample weights, so the algorithms
do not handle class imbalance identically.

MLP stopping is based on training-loss improvement, without an internal holdout split. This avoids
using labels from an internal holdout to create target encodings for the fit subset. Loss-based
stopping does not directly optimize Macro F1. Configurations remain fixed and are not changed based
on test results.

\Needspace{5\baselineskip}
\subsection{Baseline}
We fit a \code{DummyClassifier} with the \code{most\_frequent} strategy before the other algorithms. For this
dataset, the baseline always predicts class 1. On validation and test, it achieves about 57.62\%
Accuracy and a class 1 F1 of about 0.7311, but class 0 Recall is 0, Macro F1 is only about 0.3656,
and ROC-AUC is 0.5. This reference point helps identify models whose aggregate scores are driven
by predicting the majority class.

\Needspace{5\baselineskip}
\subsection{Metrics and Selection Rule}
Class 1 is designated as positive: TP is a correctly predicted high rating, FP is a high-rating
prediction for an actual score from 1 to 7, FN is a missed actual score from 8 to 10, and TN is a
correct prediction for the 1--7 group. Metrics are computed consistently with scikit-learn~\cite{metrics}:
\begin{align}
 \mathrm{Accuracy} &= \frac{TP+TN}{TP+TN+FP+FN},\\
 \mathrm{Precision}_1 &= \frac{TP}{TP+FP},\qquad
 \mathrm{Recall}_1 = \frac{TP}{TP+FN},\\
 \mathrm{Recall}_0 &= \frac{TN}{TN+FP},\qquad
 F1_k = \frac{2\,\mathrm{Precision}_k\,\mathrm{Recall}_k}
 {\mathrm{Precision}_k+\mathrm{Recall}_k},\\
 \mathrm{Macro\ F1} &= \frac{F1_0+F1_1}{2}.
\end{align}
The ROC curve plots TPR against FPR as the threshold varies; AUC is computed from continuous
probabilities or decision scores, not predicted 0/1 labels. AUC measures discrimination across
thresholds, not Accuracy or probability calibration at the selected threshold.

The project selects configurations by \textbf{validation Macro F1} to identify both classes. Ties
are broken by ROC-AUC and then Accuracy. Per-class Precision, Recall, and class 1 F1 describe
trade-offs rather than being combined into an arbitrary composite score. The notebook uses the estimators' default decision rules: a probability boundary of 0.5 for
binary LR, DT, and MLP, and a decision-score boundary of 0 for Linear SVM. These boundaries are not
tuned on the test set. For zero-denominator classification metrics, the implementation returns 0.
\FloatBarrier

\Needspace{10\baselineskip}
\section{Validation Results and Model Selection}
\Needspace{5\baselineskip}
\subsection{Comparison Across Feature Groups}
Table~\ref{tab:macro-grid} and Figure~\ref{fig:validation} compare the four algorithms using the
same metric. Full Accuracy, Precision, Recall, F1, and AUC results for all 20 configurations are
provided in Appendix~\ref{app:full}.
\begin{table}[htbp]
\centering
\caption{Validation Macro F1 for four algorithms across five input feature groups.}\label{tab:macro-grid}
\small
\begin{tabularx}{\linewidth}{l*{5}{>{\centering\arraybackslash}X}}
\toprule
Model & Initial & Metadata & Country & Count & \makecell{Target\\Encoding} \\
\midrule
LR & 0.5380 & 0.5552 & 0.5624 & 0.5638 & 0.6597 \\
DT & 0.5318 & 0.5303 & 0.5491 & 0.5855 & 0.6455 \\
SVM & 0.5380 & 0.5555 & 0.5627 & 0.5643 & 0.6594 \\
MLP & 0.3668 & 0.5232 & 0.5646 & 0.5653 & 0.6437 \\
\bottomrule
\end{tabularx}
\source{Notebook, selection summary in Section 5.6}
\end{table}
\reportfigure[1.00]{07_so_sanh_macro_f1_validation_en.pdf}{Validation Macro F1 by algorithm and feature group; the dashed line is the baseline.}{fig:validation}{Notebook, Section 5.6}
\FloatBarrier

With the Initial feature group, age, year, and missingness indicators provide limited classification
ability. LR and SVM have Macro F1 of about 0.5380 but Accuracy of about 0.5550, below the baseline's
Accuracy. This is not a contradiction: the models identify additional class 0 ratings at the cost
of some class 1 ratings. Initial MLP has class 1 Recall near 1 but class 0 Recall of only about
0.0013, yielding Macro F1 of just 0.3668. These cases illustrate why configurations should not be
selected using Accuracy or class 1 F1 alone.

Metadata increases Macro F1 and AUC for LR/SVM, but does not improve every model or metric. For
example, DT Macro F1 decreases slightly from 0.5318 to 0.5303. After adding Country, all four
algorithms improve on Metadata in this split; the gain is larger for MLP than for the other models.
These observations compare each algorithm with itself while its configuration is held fixed.

\Needspace{5\baselineskip}
\subsection{Effects of Count Features and Target Encoding}
Count features increase DT Macro F1 from 0.5491 to 0.5855. The gains are smaller for LR/SVM.
Although the Count group improves Macro F1 over Country, class 0 Recall decreases for some models;
therefore, not every type of error improves.

\begin{table}[htbp]
\centering
\caption{Logistic Regression performance as input feature groups are added sequentially.}\label{tab:lr-steps}
\small
\begin{tabularx}{\linewidth}{Yrrrr}
\toprule
Step & Accuracy & Macro F1 & $\Delta$ Macro F1 & AUC \\
\midrule
Initial $\to$ Metadata & 0.5578 & 0.5552 & +0.0172 & 0.5879 \\
Metadata $\to$ Country & 0.5653 & 0.5624 & +0.0072 & 0.5971 \\
Country $\to$ Count & 0.5670 & 0.5638 & +0.0014 & 0.6022 \\
Count $\to$ Target Encoding & 0.6610 & 0.6597 & +0.0958 & 0.7338 \\
\bottomrule
\end{tabularx}
\source{Notebook, delta tables; deltas computed before rounding}
\end{table}

Adding Target Encoding produces the clearest gains for LR/SVM among the tested groups. For LR,
Macro F1 increases by about 0.0958 over Count, while Recall increases for both classes. Because only
two ID-based rates are added at this step, the results provide predictive evidence in this split
that rating tendencies from other training records add value. They do not measure causal effects
of users or books, nor do they guarantee that the gain will persist under an ID-based or temporal split.

\begin{table}[htbp]
\centering
\caption{Validation metrics for the Target Encoding feature group.}\label{tab:te-models}
\small
\setlength{\tabcolsep}{4pt}
\begin{tabularx}{\linewidth}{l*{6}{>{\centering\arraybackslash}X}}
\toprule
 & Acc. & Prec. 1 & Recall 0 & Recall 1 & Macro F1 & AUC \\
\midrule
LR & 0.6610 & 0.7443 & 0.7071 & 0.6271 & 0.6597 & 0.7338 \\
SVM & 0.6607 & 0.7445 & 0.7080 & 0.6259 & 0.6594 & 0.7332 \\
DT & 0.6480 & 0.7207 & 0.6653 & 0.6352 & 0.6455 & 0.7154 \\
MLP & 0.6581 & 0.6875 & 0.5393 & 0.7455 & 0.6437 & 0.7116 \\
\bottomrule
\end{tabularx}
\source{Notebook, Section 5.5.4}
\end{table}

\FloatBarrier
\Needspace{8\baselineskip}
\Needspace{5\baselineskip}
\subsection{Configuration Selected by Validation}
In the supplied notebook execution, all 20 model--feature-group configurations are ranked by
validation Macro F1, ROC-AUC, and then Accuracy before the current test-evaluation cell is run.
\textbf{Logistic Regression + Target Encoding} ranks first with a Macro F1 of \textbf{0.6597},
Accuracy of \textbf{66.10\%}, and ROC-AUC of \textbf{0.7338}. Class 1 Precision is 0.7443; class 0
and class 1 Recall are 0.7071 and 0.6271, respectively.

SVM + Target Encoding is close, with a Macro F1 of about 0.6594. The difference from LR is small,
so LR can be said to lead under the selected criterion in this split only; no significance test or
multi-seed evaluation establishes a stable advantage. LR is retained for test evaluation without
refitting on the test set or reselecting the model based on test results.
\FloatBarrier

\Needspace{10\baselineskip}
\section{Test Evaluation and Generalization Scope}
\Needspace{5\baselineskip}
\subsection{Selected Model Compared with the Baseline}
The pipeline fitted on the 70\% training partition and selected by validation is evaluated unchanged.
Keeping this fitted pipeline preserves the experiment's declared protocol. A later refit on the
combined training and validation partitions would require recomputing count tables, OOF features,
and preprocessing parameters; its results would belong to a separate fitted model.
Table~\ref{tab:test} reports results for 65,051 test ratings.
\begin{table}[H]
\centering
\caption{Test results for Logistic Regression + Target Encoding and the baseline.}\label{tab:test}
\begin{tabularx}{\linewidth}{Yrr}
\toprule
Metric & Selected model & Baseline \\
\midrule
Accuracy & 0.6565 & 0.5762 \\
Class 1 Precision & 0.7407 & 0.5762 \\
Class 0 Recall & 0.7041 & 0.0000 \\
Class 0 F1 & 0.6347 & 0.0000 \\
Class 1 Recall & 0.6215 & 1.0000 \\
Class 1 F1 & 0.6758 & 0.7311 \\
Macro F1 & 0.6553 & 0.3656 \\
ROC-AUC & 0.7306 & 0.5000 \\
\bottomrule
\end{tabularx}
\source{Notebook, Section 5.7; baseline fit on training data only}
\end{table}

Accuracy increases by 8.03 percentage points, Macro F1 by 0.2897, and ROC-AUC by 0.2306 over the
baseline. Class 0 Recall increases from 0 to 0.7041, indicating that the model identifies many
ratings in the 1--7 group. In exchange, class 1 Recall is lower than for the baseline that always
predicts class 1; class 1 F1 also decreases from 0.7311 to 0.6758. Therefore, the model is not
described as outperforming the baseline on every metric.

Test Macro F1 is 0.6553, about 0.0044 below validation. This small gap in one split does not ensure
stable out-of-sample or future performance. The test set was inspected in earlier project versions,
so it is not a fully independent confirmation after development changes.

\clearpage
\subsection{Confusion Matrix and Error Types}
\reportfigure[0.76]{08_confusion_matrix_test_en.pdf}{Test confusion matrix; rows show true labels and columns show predicted labels.}{fig:confusion}{Notebook, Section 5.7}
Read Figure~\ref{fig:confusion} in the order of labels 0 and 1:
\begin{itemize}
  \item \textbf{TN = 19,412:} actual 1--7 ratings correctly predicted as class 0.
  \item \textbf{FP = 8,156:} actual 1--7 ratings incorrectly predicted as high ratings.
  \item \textbf{FN = 14,189:} actual 8--10 ratings missed and predicted as class 0.
  \item \textbf{TP = 23,294:} actual 8--10 ratings correctly identified.
\end{itemize}
Of 31,450 ratings predicted as high, 23,294 are correct, corresponding to a class 1 Precision of
about 74.07\%. The model identifies about 62.15\% of the 37,483 actual high ratings. Its class 0 F1 is
0.6347, so Macro F1 is $(0.6347+0.6758)/2\approx0.6553$. If the priority
is to find more high ratings, a threshold change could be investigated on development data; if the
priority is higher precision among high-rating predictions, the cost of FP errors should be evaluated. This
report neither specifies error costs nor adjusts the threshold, so it does not claim that the default
threshold is optimal for practical use.
\FloatBarrier

\Needspace{0.70\textheight}
\subsection{ROC and Discrimination Across Thresholds}
\reportfigure[0.76]{09_roc_test_en.pdf}{ROC curve for the final model on the test set.}{fig:roc}{Notebook, Section 5.8}
The ROC curve lies above the reference diagonal, and AUC is 0.7306, higher than the baseline value
of 0.5. This provides evidence of the model's ability to discriminate between the two classes using
continuous prediction scores rather than a single threshold. AUC does not indicate whether predicted
probabilities are well calibrated; an AUC of 0.7306 does not mean that 73.06\% of predictions are correct.

\Needspace{5\baselineskip}
\subsection{Results by Seen and Unseen IDs}
An ID is considered ``new'' if it has no ratings in the training data used to fit the model. These
segments are reported only after model selection and are not used to select the model again.
\begin{table}[htbp]
\centering
\caption{Test results by whether users and books appeared in the training set.}\label{tab:segments}
\footnotesize
\setlength{\tabcolsep}{3pt}
\begin{tabularx}{\linewidth}{Yrrrrrr}
\toprule
Segment & Ratings & Acc. & Recall 0 & Recall 1 & Macro F1 & AUC \\
\midrule
Seen user and book & 37,143 & 0.6644 & 0.6697 & 0.6609 & 0.6589 & 0.7320 \\
New user, seen book & 5,513 & 0.5674 & 0.6677 & 0.4963 & 0.5673 & 0.6118 \\
Seen user, new book & 19,908 & 0.6832 & 0.7390 & 0.6361 & 0.6832 & 0.7632 \\
New user and book & 2,487 & 0.5219 & 0.9260 & 0.1097 & 0.4234 & 0.5935 \\
\bottomrule
\end{tabularx}
\source{Notebook, Section 5.7; the four segments sum to 65,051 ratings}
\end{table}

The seen-user/seen-book segment contains 37,143 ratings and has a Macro F1 of 0.6589. When both IDs
are new, Macro F1 falls to 0.4234 and class 1 Recall is 0.1097. In this segment, both target-encoding
features fall back to the training prior and both count features are 0, leaving only metadata,
age, country, and missingness information to distinguish records. This
result cautions against generalizing aggregate performance to all new-user/new-book cases.

The seen-user/new-book segment achieves a Macro F1 of 0.6832, higher than the segment where both
are seen. The segments differ in composition, class distribution, and size, so this does not imply
that ``new books improve the model.'' The current analysis is descriptive; general conclusions
require an appropriate ID-based split and repeated evaluation.
\FloatBarrier

\Needspace{10\baselineskip}
\section{Discussion, Applications, and Limitations}
\Needspace{5\baselineskip}
\subsection{Interpretation of Findings}
The EDA shows substantial missingness in age, an uneven distribution of rating counts, and a
difference between popularity and mean score. Basic demographic and metadata features do not
provide strong discrimination. In these experiments, ID-based rating-tendency statistics produce
the clearest gains for LR/SVM. This is predictive value under the current design, not evidence that
any individual feature causes rating behavior.

For applications, the model could support estimates of high-rating likelihood when user/book
information and other ratings are available for lookup. Precision and Recall help inform trade-offs
for a future objective. However, before using the model to rank recommendations, a ranking-evaluation
design and data on candidate books that have not been rated are needed. The current metrics do not
establish purchasing benefit, revenue, or business impact.

\clearpage
\subsection{Limitations to Consider}
\begin{enumerate}
  \item \textbf{Missingness and data quality:} age, publication year, and book metadata have
  substantial missingness. Imputation and missingness indicators do not recover the original
  information. Age/year thresholds are conventions that may mark some valid values as missing;
  country parsed from Location may remain incorrect or inconsistent.
  \item \textbf{Rating-level split:} the same users and books can occur across multiple splits.
  The data have no timestamps to establish that all statistics would be available before a rating
  in a real deployment. The results are not a pure cold-start or future-prediction evaluation.
  \item \textbf{Label scope:} excluding zero-valued ratings matches the explicit-rating task
  definition but limits conclusions to records scored from 1 to 10. This group may differ from the
  population of users or books that have not been rated.
  \item \textbf{Statistical inference:} EDA uses the full sample, and ratings may be dependent by
  user or book. User-level rates have different precision because users contribute different
  numbers of ratings; small $p$-values are not used to claim causality or stable population-level
  differences.
  \item \textbf{Selection and independence:} all 20 configurations share one validation split and
  one seed, and hyperparameters are not exhaustively optimized. The test set was inspected in earlier
  project versions, so the current evaluation is not a fully independent post-development confirmation.
  \item \textbf{Distribution differences:} training uses OOF rates and counts with the current
  rating removed, while validation/test use statistics from all training data. Smoothing reduces
  but does not eliminate these differences. MLP stops based on training loss rather than directly
  optimizing Macro F1, and class balancing differs across algorithms.
\end{enumerate}

\clearpage
\section{Conclusion and Future Work}
\Needspace{5\baselineskip}
\subsection{Conclusion}
The project covers the full workflow from relational data validation, cleaning, EDA, and label
definition to feature engineering, training, and evaluation. The classification dataset contains
433,671 ratings. Retaining ISBNs as strings, normalizing keys, preserving ratings without metadata,
and fitting preprocessing on training data are important for consistency.

Among the 20 evaluated configurations, \textbf{Logistic Regression + Target Encoding} is selected
by validation Macro F1. On test, the model achieves \textbf{65.65\% Accuracy, 0.6553 Macro F1, and
0.7306 ROC-AUC}. Compared with the baseline, it improves Accuracy and the equal-weighted summary of both classes,
while class 1 Recall and F1 decrease. Thus the gain in Macro F1 reflects recovery of class 0 rather
than an improvement on every measure. Weaker results when both IDs are new show why performance should be considered
by use case rather than summarized by one aggregate score.

\Needspace{5\baselineskip}
\subsection{Next Steps}
\begin{enumerate}
  \item \textbf{Validate on new data:} prepare a holdout set that has never been inspected and
  evaluate across additional seeds or with cross-validation. Within every fold, refit preprocessing,
  count features, and OOF encodings using only that fold's fit partition.
  \item \textbf{Choose a suitable generalization design:} try user-, book-, or time-based splits if
  timestamps become available. Evaluate unseen-ID cases and different history densities separately.
  \item \textbf{Optimize for the intended objective:} select hyperparameters and thresholds on
  development data, and assess PR-AUC, probability calibration, FP/FN costs, and result variability.
  \item \textbf{Improve representations:} normalize country values more carefully and consider book
  title text, genres, or descriptions if suitable sources are available. Evaluate additions under a
  consistent design.
  \item \textbf{If moving to recommendations:} construct a candidate-book set and evaluate ranking
  appropriately; do not treat the current classification results as evidence of recommendation quality.
\end{enumerate}

\clearpage
\Needspace{10\baselineskip}
\section*{Reproducibility and Data Provenance}
\addcontentsline{toc}{section}{Reproducibility and Data Provenance}
The supplied \nolinkurl{mid_term_code.ipynb} contains 70 cells, including 41 code cells with
execution counts 1--41 and no saved error or warning outputs. The notebook was rerun from
beginning to end in a fresh kernel, including all 20 model--feature-group configurations.
Its saved tables are the source for the comparisons in this report. All 140 metric values in
the complete validation table match the notebook to four decimal places; the test metrics,
confusion-matrix counts, and ID-segment results agree as well. This verifies consistency with
the supplied notebook; rerunning the same split does not create an independent holdout.

\begin{table}[htbp]
\centering
\caption{Environment recorded in the notebook run used for this report.}\label{tab:environment}
\begin{tabular}{ll}
\toprule
Component & Version \\
\midrule
Python & 3.14.8 \\
numpy & 2.5.4 \\
pandas & 3.1.0 \\
matplotlib & 3.11.2 \\
scipy & 1.18.1 \\
scikit-learn & 1.9.1 \\
\bottomrule
\end{tabular}
\end{table}

The data-split seed is 56; the algorithm and OOF seeds are 42. To reproduce the analysis, place the
three CSV files in \nolinkurl{DATA_SET/}, install numpy, pandas, matplotlib, scipy, scikit-learn,
and Jupyter, open the notebook from the project root, and run all cells in a fresh kernel. The CSV files must
be under \nolinkurl{DATA_SET/}, matching the notebook's \code{DATA\_DIR} setting.
Results may differ slightly across library versions or numerical backends.

\noindent\textbf{Notebook version fingerprint:} SHA-256 begins with
\nolinkurl{27eb7d78452c15fc}. Table~\ref{tab:environment} records the environment printed by
that notebook execution. The fingerprint identifies the exact notebook file supplied with this report.

All nine figures have been regenerated with English labels. Descriptive charts use the supplied
data; the validation chart uses the saved 20-configuration table; the confusion matrix and ROC use
the selected LR predictions whose reported results match the notebook. Vector PDF charts are included in
the source package. Tables use a period as the decimal separator and show metrics to four decimal
places. Deltas use unrounded notebook values, so the final digit can differ from subtracting two
rounded entries.

\clearpage
\begin{thebibliography}{9}
\addcontentsline{toc}{section}{References}
\bibitem{dataset}
Möbius. \emph{Book Recommendation Dataset}. Kaggle.
\url{https://www.kaggle.com/datasets/arashnic/book-recommendation-dataset}.
\bibitem{description}
Möbius. \emph{Practice Recommendation Algorithms with Book Recommendation Dataset}.
Description of Users, Books, Ratings, and the meaning of rating scores.
\url{https://www.kaggle.com/discussions/general/202108}.
\bibitem{metrics}
Scikit-learn developers. \emph{Metrics and scoring: quantifying the quality of predictions},
section on classification metrics.
\url{https://scikit-learn.org/stable/modules/model_evaluation.html}.
\bibitem{pitfalls}
Scikit-learn developers. \emph{Common pitfalls and recommended practices}, section on data leakage.
\url{https://scikit-learn.org/stable/common_pitfalls.html}.
\bibitem{encoding}
Scikit-learn developers. \emph{TargetEncoder}, documentation on cross-fitting, smoothing,
and unseen categories.
\url{https://scikit-learn.org/stable/modules/generated/sklearn.preprocessing.TargetEncoder.html}.
\bibitem{notebook}
\emph{mid\_term\_code.ipynb}. Project source code and experiment outputs.
\url{https://github.com/Bluenip/Data_analyst_python_midterm_project}.
\end{thebibliography}

\clearpage
\appendix
\Needspace{10\baselineskip}
\section{Complete Validation Results for 20 Configurations}\label{app:full}
LR: Logistic Regression; DT: Decision Tree; SVM: Linear SVM; MLP: multilayer perceptron.
Acc.: Accuracy; P1: class 1 Precision; R0/R1: class 0/1 Recall; F1$_1$: class 1 F1;
MF1: Macro F1; AUC: ROC-AUC. All 20 configurations use the same validation set; this table does
not include test results.

\begingroup
\footnotesize
\setlength{\tabcolsep}{3pt}
\renewcommand{\arraystretch}{1.2}
\begin{longtable}{L{2.4cm}lrrrrrrr}
\caption{All validation metrics for five feature groups and four algorithms.}\label{tab:all-results}\\
\toprule
Group & Model & Acc. & P1 & R0 & R1 & F1$_1$ & MF1 & AUC \\
\midrule
\endfirsthead
\multicolumn{9}{l}{\textit{Table \thetable{} (continued)}}\\
\toprule
Group & Model & Acc. & P1 & R0 & R1 & F1$_1$ & MF1 & AUC \\
\midrule
\endhead
\midrule
\multicolumn{9}{r}{\textit{Continued on next page}}\\
\endfoot
\bottomrule
\endlastfoot
Initial & LR & 0.5550 & 0.6066 & 0.4286 & 0.6479 & 0.6266 & 0.5380 & 0.5455 \\
Initial & DT & 0.5319 & 0.6229 & 0.6085 & 0.4755 & 0.5393 & 0.5318 & 0.5586 \\
Initial & SVM & 0.5550 & 0.6066 & 0.4286 & 0.6479 & 0.6266 & 0.5380 & 0.5455 \\
Initial & MLP & 0.5764 & 0.5764 & 0.0013 & 0.9993 & 0.7311 & 0.3668 & 0.5516 \\
Metadata & LR & 0.5578 & 0.6341 & 0.5684 & 0.5500 & 0.5891 & 0.5552 & 0.5879 \\
Metadata & DT & 0.5318 & 0.6469 & 0.6937 & 0.4126 & 0.5038 & 0.5303 & 0.5793 \\
Metadata & SVM & 0.5580 & 0.6345 & 0.5695 & 0.5496 & 0.5890 & 0.5555 & 0.5879 \\
Metadata & MLP & 0.5896 & 0.6040 & 0.2555 & 0.8353 & 0.7011 & 0.5232 & 0.5976 \\
Country & LR & 0.5653 & 0.6402 & 0.5717 & 0.5606 & 0.5978 & 0.5624 & 0.5971 \\
Country & DT & 0.5491 & 0.6475 & 0.6465 & 0.4775 & 0.5496 & 0.5491 & 0.5920 \\
Country & SVM & 0.5654 & 0.6407 & 0.5733 & 0.5596 & 0.5974 & 0.5627 & 0.5971 \\
Country & MLP & 0.5928 & 0.6246 & 0.3992 & 0.7352 & 0.6754 & 0.5646 & 0.6082 \\
Count & LR & 0.5670 & 0.6408 & 0.5690 & 0.5655 & 0.6008 & 0.5638 & 0.6022 \\
Count & DT & 0.5882 & 0.6630 & 0.5990 & 0.5802 & 0.6189 & 0.5855 & 0.6324 \\
Count & SVM & 0.5674 & 0.6414 & 0.5703 & 0.5653 & 0.6009 & 0.5643 & 0.6022 \\
Count & MLP & 0.6056 & 0.6248 & 0.3554 & 0.7896 & 0.6976 & 0.5653 & 0.6333 \\
Target Encoding & LR & 0.6610 & 0.7443 & 0.7071 & 0.6271 & 0.6807 & 0.6597 & 0.7338 \\
Target Encoding & DT & 0.6480 & 0.7207 & 0.6653 & 0.6352 & 0.6753 & 0.6455 & 0.7154 \\
Target Encoding & SVM & 0.6607 & 0.7445 & 0.7080 & 0.6259 & 0.6801 & 0.6594 & 0.7332 \\
Target Encoding & MLP & 0.6581 & 0.6875 & 0.5393 & 0.7455 & 0.7153 & 0.6437 & 0.7116 \\
\end{longtable}
\endgroup

\noindent\textit{Note:} Feature groups are ordered by the sequence in which inputs are added, not by
final rank. The selected configuration has the highest MF1; ties are broken by AUC and then Acc.
Each algorithm's configuration is held fixed across groups, but the project does not establish that
any configuration is optimal across all representations, hyperparameters, or data splits.

\end{document}
